\section{Regularidad dinámica del proceso}\label{sec:regularidad}

La comparación entrópica ya permite transferir eventos suficientemente raros.
Para aplicarla a las reparaciones, primero necesitamos saber con qué frecuencia
se eligen las líneas que ocuparían las dos reinas de una reparación. Esas
probabilidades dependen de los pesos y de los grados disponibles en cada instante.

\subsection{Pesos reescalados y balance inicial}

Para $1\le h\le K$ reescalamos sólo los pesos reales por el factor común $x_*^{-2}$, lo que no
cambia la ley: $v_e=\varpi(e)/x_*^2=r_Rr_Cu_Du_A$. Para $h=0$, $v_e$ es directamente el peso real
de la Definición~\ref{def:referencia_cero}. Por el Teorema~\ref{thm:perfil}(a), $c\le v_e\le C$ con
$c=\underline r^2\kappa^2$ y $C=M^2$. Ponemos
\[
 W_0=\sum_ev_e\asymp n^2,\quad \deg_v^0=\sum_{e\ni v}v_e,\quad \pi_e=\frac{qv_e}{W_0},\quad
 p_v=\sum_{e\ni v}\pi_e,\quad U_f^0=\sum_{v\in f}u_v .
\]
Los pesos virtuales $u_v$ se conservan; los ficticios tienen $u_v=1$ y $p_v=0$.

\begin{lema}[Balance inicial]\label{lem:balance}
Uniformemente en $0\le h\le K$, para todo vértice físico $v$ y toda arista real $e$,
\begin{equation}\label{eq:balance}
 p_v+u_v=1+O(n^{-1}),\qquad U_f^0=H+O(1),\qquad p_v\le C,\qquad \pi_e\le C/n .
\end{equation}
Además $\sum_{v\in f}p_v=q$ en cada familia y, para diagonales y antidiagonales,
$p_v\le1-\eta$ con $\eta=\kappa/2$ y $n$ grande.
\end{lema}

\begin{proof}
Sea $\rho_n=q/n$. Por cuadratura, $W_0=n^2\rho_n/b\,(1+O(n^{-1}))$. En una fila $r$, por la
ecuación de Euler~\eqref{eq:euler},
$q\deg_v^0/W_0=b\,r_R(\bar r)J_R(\bar r)+O(n^{-1})=1-u_R+O(n^{-1})$; en una diagonal,
$q\deg_v^0/W_0=b\,u_DJ_D+O(n^{-1})=1-u_D+O(n^{-1})$; columnas y antidiagonales son análogas.
Las cotas y cuadraturas son uniformes hasta el borde del perfil objetivo. En cada familia cada
arista real aparece exactamente una vez, luego $\sum_{v\in f}\deg_v^0=W_0$ y
$\sum_{v\in f}p_v=q$; sumando el balance puntual, $U_f^0=N-q+O(N/n)=H+O(1)$, incluidos los
ficticios. En diagonales $u_v\ge\kappa$, de donde $p_v\le1-\kappa+O(n^{-1})\le1-\eta$. En $h=0$
el balance es $p_v=1+O(n^{-1})$, $u_v=0$ en filas y columnas físicas, la suma virtual de cada una
de esas familias es $H$ (sólo ficticios) y las diagonales no cambian.
\end{proof}

\emph{Relojes.} $t$ cuenta todas las piezas colocadas, $z_t=(N-t)/N$ es la fracción activa y
$q_t$ el número de piezas reales pendientes. Ni $z_t$ ni $q_t/n$ sustituyen al parámetro fijo
$x_*$ en los pesos.
